\documentclass[aspectratio=169,11pt]{beamer}
\usetheme{Madrid}
\usecolortheme{dolphin}
\setbeamertemplate{navigation symbols}{}
\setbeamertemplate{caption}[numbered]

\usepackage[utf8]{inputenc}
\usepackage[T1]{fontenc}
\usepackage[spanish,es-noshorthands]{babel}
\usepackage{amsmath,amssymb,amsthm,mathtools}
\usepackage{tikz}
\usepackage{pgfplots}
\pgfplotsset{compat=1.18}
\usepackage{booktabs}
\usepackage{array}

\title[Prob. y Estadística]{Clase 27: Pruebas de hipótesis para dos muestras}
\subtitle{Unidad 7: Prueba de hipótesis}
\institute[UNS]{Universidad Nacional del Sur \\ Departamento de Matemática}
\author{Probabilidad y Estadística (5765)}
\date{}

\begin{document}

\begin{frame}
\titlepage
\end{frame}

\begin{frame}{Contenidos}
\tableofcontents
\end{frame}

\section{Introducción}

\begin{frame}{Comparación de dos poblaciones}
Es habitual comparar dos poblaciones a través de sus medias o sus varianzas: dos procesos de producción, dos proveedores, un tratamiento vs. un control, mediciones antes y después de una intervención, etc.

\begin{block}{Situaciones a distinguir}
\begin{itemize}
\item \textbf{Muestras independientes} con varianzas poblacionales \textbf{conocidas}.
\item \textbf{Muestras independientes} con varianzas poblacionales \textbf{desconocidas pero iguales} (se estima una varianza combinada, \emph{pooled}).
\item \textbf{Muestras apareadas} (pareadas o dependientes): cada observación de una muestra está naturalmente asociada a una de la otra (mismo individuo antes/después, mismo lote medido por dos métodos, etc.).
\item Comparación de \textbf{varianzas} de dos poblaciones (prueba $F$).
\end{itemize}
\end{block}
\end{frame}

\begin{frame}{¿Muestras independientes o apareadas?}
\begin{alertblock}{Pregunta clave del diseño}
Antes de elegir el procedimiento, hay que identificar si existe una \textbf{correspondencia natural} entre cada observación de una muestra y una de la otra.
\end{alertblock}

\begin{columns}[t]
\begin{column}{0.5\textwidth}
\textbf{Independientes}
\begin{itemize}
\item Dos máquinas distintas.
\item Dos proveedores distintos.
\item Grupo control vs. grupo tratado (sujetos distintos).
\end{itemize}
\end{column}
\begin{column}{0.5\textwidth}
\textbf{Apareadas}
\begin{itemize}
\item Mismo sujeto antes/después.
\item Misma pieza medida por dos instrumentos.
\item Gemelos, o unidades emparejadas por un criterio.
\end{itemize}
\end{column}
\end{columns}
\vspace{0.3cm}
Confundir el diseño lleva a elegir un estadístico incorrecto y a conclusiones erróneas.
\end{frame}

\section{Diferencia de medias: varianzas conocidas}

\begin{frame}{Prueba para $\mu_1-\mu_2$, varianzas conocidas}
Sean $X_1,\dots,X_{n_1}$ muestra de $N(\mu_1,\sigma_1^2)$ y $Y_1,\dots,Y_{n_2}$ muestra de $N(\mu_2,\sigma_2^2)$, \textbf{independientes}, con $\sigma_1^2,\sigma_2^2$ \textbf{conocidas}. Se desea probar
\[
H_0:\mu_1-\mu_2=\delta_0 \qquad (\text{usualmente } \delta_0=0).
\]

\begin{block}{Estadístico de prueba}
\[
Z=\frac{(\overline{X}-\overline{Y})-\delta_0}{\sqrt{\dfrac{\sigma_1^2}{n_1}+\dfrac{\sigma_2^2}{n_2}}} \ \overset{H_0}{\sim}\ N(0,1).
\]
\end{block}

Las regiones críticas son análogas al caso de una media: $|Z|>z_{\alpha/2}$ (bilateral), $Z>z_\alpha$ o $Z<-z_\alpha$ (unilaterales), según $H_1$.
\end{frame}

\begin{frame}{Ejemplo: comparación de dos máquinas, varianzas conocidas}
\begin{example}
Dos máquinas envasadoras de un mismo producto tienen desviaciones estándar históricas conocidas $\sigma_1=5$ g y $\sigma_2=4$ g. Se toman $n_1=40$ envases de la máquina 1 ($\overline{x}=252$ g) y $n_2=35$ de la máquina 2 ($\overline{y}=249$ g). Con $\alpha=0.05$, ¿difiere el peso medio entregado por ambas máquinas?

\textbf{1. Hipótesis:} $H_0:\mu_1-\mu_2=0$ vs. $H_1:\mu_1-\mu_2\neq0$.

\textbf{2. Estadístico:} $Z=\dfrac{(\overline{X}-\overline{Y})-0}{\sqrt{\sigma_1^2/n_1+\sigma_2^2/n_2}}$.

\textbf{3. Región crítica:} $|Z|>z_{0.025}=1.96$.

\textbf{4. Cálculo:}
\[
Z=\frac{(252-249)-0}{\sqrt{25/40+16/35}}=\frac{3}{\sqrt{0.625+0.457}}=\frac{3}{\sqrt{1.082}}=\frac{3}{1.040}=2.88.
\]

\textbf{5. Decisión:} $2.88>1.96 \Rightarrow$ se \textbf{rechaza} $H_0$.

\textbf{6. Conclusión:} hay evidencia significativa de que el peso medio entregado por ambas máquinas difiere.
\end{example}
\end{frame}

\begin{frame}{El ejemplo anterior con valor p}
\begin{example}
Con $z=2.88$ (prueba bilateral), el valor p es
\[
p=2\cdot P(Z>2.88)=2\times0.0020=0.0040.
\]
Como $p=0.0040<\alpha=0.05$ (y también $<0.01$), la evidencia en contra de $H_0$ es muy fuerte: la diferencia observada entre las dos máquinas difícilmente se debe al azar si realmente $\mu_1=\mu_2$.
\end{example}

\begin{block}{Intervalo de confianza asociado}
Un IC del $95\%$ para $\mu_1-\mu_2$ es
\[
(\overline{x}-\overline{y}) \pm z_{0.025}\sqrt{\tfrac{\sigma_1^2}{n_1}+\tfrac{\sigma_2^2}{n_2}} = 3 \pm 1.96\times1.040 = 3\pm2.04 = (0.96;\,5.04),
\]
que no contiene al $0$, consistente con el rechazo de $H_0:\mu_1-\mu_2=0$.
\end{block}
\end{frame}

\section{Diferencia de medias: varianzas desconocidas e iguales}

\begin{frame}{Prueba $t$ agrupada (\emph{pooled})}
Cuando $\sigma_1^2$ y $\sigma_2^2$ son desconocidas pero puede suponerse $\sigma_1^2=\sigma_2^2=\sigma^2$ (poblaciones normales, muestras independientes), se estima $\sigma^2$ combinando ambas cuasivarianzas muestrales:

\begin{block}{Varianza combinada (\emph{pooled})}
\[
S_p^2 = \frac{(n_1-1)S_1^2+(n_2-1)S_2^2}{n_1+n_2-2}.
\]
\end{block}

\begin{block}{Estadístico de prueba}
\[
T=\frac{(\overline{X}-\overline{Y})-\delta_0}{S_p\sqrt{\dfrac{1}{n_1}+\dfrac{1}{n_2}}} \ \overset{H_0}{\sim}\ t_{n_1+n_2-2}.
\]
\end{block}

Regiones críticas análogas, usando $t_{n_1+n_2-2,\,\gamma}$ en lugar de $z_\gamma$.
\end{frame}

\begin{frame}{Ejemplo: comparación de dos proveedores}
\begin{example}
Se compara la resistencia a la rotura (en kg) de cables de dos proveedores, suponiendo varianzas poblacionales iguales. Proveedor A: $n_1=12$, $\overline{x}=436$, $s_1^2=98$. Proveedor B: $n_2=10$, $\overline{y}=428$, $s_2^2=110$. Con $\alpha=0.05$, ¿el proveedor A ofrece mayor resistencia media?

\textbf{1. Hipótesis:} $H_0:\mu_1-\mu_2\le0$ vs. $H_1:\mu_1-\mu_2>0$.

\textbf{2. Varianza combinada:}
\[
S_p^2=\frac{11(98)+9(110)}{20}=\frac{1078+990}{20}=\frac{2068}{20}=103.4, \quad S_p=10.17.
\]

\textbf{3. Región crítica} (gl $=20$): $T>t_{20,\,0.05}=1.725$.

\textbf{4. Cálculo:}
\[
t=\frac{(436-428)-0}{10.17\sqrt{\tfrac{1}{12}+\tfrac{1}{10}}}=\frac{8}{10.17\times0.4282}=\frac{8}{4.354}=1.84.
\]

\textbf{5. Decisión:} $1.84>1.725 \Rightarrow$ se \textbf{rechaza} $H_0$.

\textbf{6. Conclusión:} con un $5\%$ de significación, hay evidencia de que el proveedor A entrega cables con mayor resistencia media.
\end{example}
\end{frame}

\begin{frame}{Observación: varianzas desconocidas y distintas}
\begin{alertblock}{Cuando no puede suponerse $\sigma_1^2=\sigma_2^2$}
Si las varianzas poblacionales son desconocidas y \textbf{no} puede suponerse que sean iguales, el estadístico $t$ agrupado no es válido. Existe una extensión (aproximación de Welch–Satterthwaite) que utiliza un estadístico $t$ con grados de libertad ajustados, distinta de la fórmula vista. Su desarrollo excede el alcance de esta materia; aquí siempre trabajaremos bajo el supuesto de \textbf{varianzas iguales} cuando estas son desconocidas.
\end{alertblock}
\begin{block}{Cómo decidir si es razonable suponer $\sigma_1^2=\sigma_2^2$}
Este supuesto puede evaluarse formalmente con la prueba $F$ de igualdad de varianzas, que se presenta más adelante en esta misma clase.
\end{block}
\end{frame}

\section{Muestras apareadas}

\begin{frame}{Prueba para datos apareados}
Cuando las dos muestras \textbf{no son independientes} (mismo sujeto/unidad medido dos veces, o pares naturalmente asociados), se trabaja con las \textbf{diferencias} $D_i=X_i-Y_i$, $i=1,\dots,n$, reduciendo el problema a una prueba de una sola media.

\begin{block}{Estadístico de prueba}
Si $D_1,\dots,D_n$ son aproximadamente $N(\mu_D,\sigma_D^2)$ con $\sigma_D^2$ desconocida:
\[
T=\frac{\overline{D}-\delta_0}{S_D/\sqrt{n}} \ \overset{H_0}{\sim}\ t_{n-1}, \qquad H_0:\mu_D=\delta_0 \ (\text{usualmente } 0).
\]
donde $\overline{D}$ y $S_D$ son la media y el desvío muestral de las diferencias.
\end{block}

\begin{alertblock}{Por qué se usa este enfoque}
Al trabajar con las diferencias se elimina la variabilidad debida a diferencias entre unidades (cada unidad actúa como su propio control), lo que suele aumentar la potencia de la prueba respecto de tratar las muestras como independientes.
\end{alertblock}
\end{frame}

\begin{frame}{Ejemplo: prueba apareada (antes/después)}
\begin{example}
Se evalúa un programa de entrenamiento midiendo el tiempo (en segundos) que tardan $8$ operarios en completar una tarea, antes y después del entrenamiento:
\begin{center}
\small
\begin{tabular}{@{}lcccccccc@{}}
\toprule
Operario & 1 & 2 & 3 & 4 & 5 & 6 & 7 & 8 \\
\midrule
Antes & 25 & 28 & 22 & 30 & 27 & 24 & 26 & 29 \\
Después & 23 & 24 & 22 & 27 & 25 & 24 & 22 & 25 \\
Diferencia $d_i$ & 2 & 4 & 0 & 3 & 2 & 0 & 4 & 4 \\
\bottomrule
\end{tabular}
\end{center}
Se obtiene $\overline{d}=2.375$ y $s_D=1.685$. Con $\alpha=0.05$, ¿el entrenamiento reduce el tiempo medio de ejecución?

\textbf{Hipótesis:} $H_0:\mu_D\le0$ vs. $H_1:\mu_D>0$ (con $D=$Antes$-$Después).

\textbf{Región crítica} (gl$=7$): $T>t_{7,\,0.05}=1.895$.

\textbf{Cálculo:} $\displaystyle t=\frac{2.375-0}{1.685/\sqrt{8}}=\frac{2.375}{0.5958}=3.99$.

\textbf{Decisión y conclusión:} $3.99>1.895\Rightarrow$ se rechaza $H_0$: el entrenamiento reduce significativamente el tiempo medio de ejecución.
\end{example}
\end{frame}

\section{Igualdad de dos varianzas}

\begin{frame}{Prueba $F$ para igualdad de varianzas}
Sean $X_1,\dots,X_{n_1}\sim N(\mu_1,\sigma_1^2)$ y $Y_1,\dots,Y_{n_2}\sim N(\mu_2,\sigma_2^2)$ muestras independientes. Se desea probar
\[
H_0:\sigma_1^2=\sigma_2^2.
\]

\begin{block}{Estadístico de prueba}
\[
F=\frac{S_1^2}{S_2^2} \ \overset{H_0}{\sim}\ F_{n_1-1,\,n_2-1},
\]
la distribución $F$ de Snedecor con $n_1-1$ grados de libertad en el numerador y $n_2-1$ en el denominador.
\end{block}

\begin{table}[]
\centering
\small
\begin{tabular}{@{}lc@{}}
\toprule
$H_1$ & Región crítica \\
\midrule
$\sigma_1^2\neq\sigma_2^2$ & $F<F_{n_1-1,n_2-1,\,1-\alpha/2}$ o $F>F_{n_1-1,n_2-1,\,\alpha/2}$ \\
$\sigma_1^2>\sigma_2^2$ & $F>F_{n_1-1,n_2-1,\,\alpha}$ \\
\bottomrule
\end{tabular}
\end{table}
\begin{alertblock}{Convención práctica}
Para el caso bilateral es usual colocar en el numerador la muestra con mayor $s^2$, de modo que $F\ge1$, y comparar solo con el valor crítico superior $F_{n_1-1,n_2-1,\alpha/2}$.
\end{alertblock}
\end{frame}

\begin{frame}{Forma de la distribución $F$}
\centering
\begin{tikzpicture}
\begin{axis}[
  width=9.5cm, height=5cm,
  axis lines=left, ytick=\empty,
  xlabel={$F$}, ylabel={densidad},
  xmin=0, xmax=5, ymin=0,
  legend pos=north east, legend style={font=\footnotesize},
]
\addplot[domain=0.05:5, samples=150, thick, blue] {(x^1)/((1+2*x)^3)};
\addlegendentry{$F_{4,10}$}
\addplot[domain=0.05:5, samples=150, thick, red, dashed] {(x^5)/((1+1.2*x)^8)*3};
\addlegendentry{$F_{12,20}$ (más concentrada en $1$)}
\end{axis}
\end{tikzpicture}
\vspace{0.2cm}

Al igual que $\chi^2$, la distribución $F$ toma solo valores positivos y es asimétrica; su forma depende de \textbf{dos} grados de libertad (numerador y denominador).
\end{frame}

\begin{frame}{Ejemplo: comparación unilateral de variabilidad}
\begin{example}
Un nuevo instrumento de medición se propone como más preciso (menor variabilidad) que el instrumento estándar. Instrumento estándar: $n_1=16$, $s_1^2=2.8$. Instrumento nuevo: $n_2=13$, $s_2^2=1.1$. Con $\alpha=0.05$, ¿hay evidencia de que el nuevo instrumento es más preciso?

\textbf{1. Hipótesis:} $H_0:\sigma_1^2\le\sigma_2^2$ vs. $H_1:\sigma_1^2>\sigma_2^2$ (equivalente a $H_0:\sigma_1^2=\sigma_2^2$).

\textbf{2. Estadístico:} $\displaystyle F=\frac{S_1^2}{S_2^2}\sim F_{15,\,12}$ bajo $H_0$.

\textbf{3. Región crítica:} $F>F_{15,\,12,\,0.05}=2.62$.

\textbf{4. Cálculo:} $\displaystyle F=\frac{2.8}{1.1}=2.55$.

\textbf{5. Decisión:} $2.55<2.62\Rightarrow$ \textbf{no se rechaza} $H_0$.

\textbf{6. Conclusión:} al $5\%$, no hay evidencia estadísticamente suficiente de que el nuevo instrumento sea más preciso, aunque el resultado es sugestivo y amerita ampliar la muestra.
\end{example}
\end{frame}

\begin{frame}{Ejemplo: comparación de variabilidad de dos procesos}
\begin{example}
Se comparan dos líneas de producción respecto de la variabilidad en el diámetro (mm) de una pieza. Línea A: $n_1=13$, $s_1^2=0.048$. Línea B: $n_2=10$, $s_2^2=0.019$. Con $\alpha=0.10$, ¿hay evidencia de que las variabilidades difieren?

\textbf{1. Hipótesis:} $H_0:\sigma_1^2=\sigma_2^2$ vs. $H_1:\sigma_1^2\neq\sigma_2^2$.

\textbf{2. Estadístico} (numerador = mayor varianza): $\displaystyle F=\frac{S_1^2}{S_2^2}\sim F_{12,\,9}$ bajo $H_0$.

\textbf{3. Región crítica:} $F>F_{12,\,9,\,0.05}=3.07$.

\textbf{4. Cálculo:}
\[
F=\frac{0.048}{0.019}=2.53.
\]

\textbf{5. Decisión:} $2.53<3.07 \Rightarrow$ \textbf{no se rechaza} $H_0$.

\textbf{6. Conclusión:} al nivel $10\%$, no hay evidencia suficiente de que las variabilidades de ambas líneas difieran.
\end{example}
\end{frame}

\section{Resumen}

\begin{frame}{Resumen: pruebas para dos poblaciones}
\begin{table}[]
\centering
\scriptsize
\begin{tabular}{@{}p{3.0cm}p{3.2cm}p{2.9cm}p{2.6cm}@{}}
\toprule
Situación & Estadístico & Distribución & Supuestos \\
\midrule
Medias, $\sigma_1^2,\sigma_2^2$ conocidas & $Z=\dfrac{(\overline{X}-\overline{Y})-\delta_0}{\sqrt{\sigma_1^2/n_1+\sigma_2^2/n_2}}$ & $N(0,1)$ & independencia \\[1.4em]
Medias, varianzas desc. iguales & $T=\dfrac{(\overline{X}-\overline{Y})-\delta_0}{S_p\sqrt{1/n_1+1/n_2}}$ & $t_{n_1+n_2-2}$ & independencia, $\sigma_1^2=\sigma_2^2$ \\[1.4em]
Muestras apareadas & $T=\dfrac{\overline{D}-\delta_0}{S_D/\sqrt{n}}$ & $t_{n-1}$ & normalidad de $D_i$ \\[1.2em]
Igualdad de varianzas & $F=\dfrac{S_1^2}{S_2^2}$ & $F_{n_1-1,n_2-1}$ & independencia, normalidad \\
\bottomrule
\end{tabular}
\end{table}

\begin{itemize}
\item La elección entre \emph{pooled t} y datos apareados depende del \textbf{diseño del estudio}, no es opcional.
\item Próxima clase: pruebas de hipótesis basadas en la distribución $\chi^2$ para bondad de ajuste, independencia y homogeneidad.
\end{itemize}
\end{frame}

\end{document}
